% ---------- Fonctionnement général ----------
\begin{frame}{Fonctionnement général du parcours chirurgical}

  \begin{columns}[T]

    \begin{column}{0.48\textwidth}

      \begin{block}{Planification des interventions}
        Une intervention doit être programmée en tenant compte :
        \begin{itemize}
          \item de la priorité du patient ;
          \item de la disponibilité du chirurgien ;
          \item des blocs opératoires ;
          \item des équipes d'anesthésie et de soins ;
          \item des lits disponibles après l'intervention.
        \end{itemize}
      \end{block}

    \end{column}

    \begin{column}{0.48\textwidth}

      \begin{block}{Hospitalisation avant / après intervention}
        \textbf{Avant :}
        \begin{itemize}
          \item admission du patient ;
          \item préparation préopératoire ;
          \item passage éventuel par le SAS de pré-anesthésie ;
          \item entrée au bloc opératoire.
        \end{itemize}

        \textbf{Après :}
        \begin{itemize}
          \item passage en salle de réveil (SSPI) ;
          \item hospitalisation ambulatoire ou conventionnelle ;
          \item sortie après récupération.
        \end{itemize}
      \end{block}

    \end{column}

  \end{columns}

\end{frame}

% ---------- Contraintes et enjeux ----------
\begin{frame}{Contraintes et enjeux}

  \begin{columns}[T]

    \begin{column}{0.48\textwidth}

      \begin{block}{Contraintes}
        \begin{itemize}
          \item Capacité limitée : \textbf{45 lits} et un nombre limité de blocs.
          \item Durées variables des interventions et des séjours.
          \item Intégration des urgences dans le planning.
          \item Contraintes des chirurgiens (vacations et salles spécifiques).
        \end{itemize}
      \end{block}

    \end{column}

    \begin{column}{0.48\textwidth}

      \begin{block}{Enjeux}
        \begin{itemize}
          \item Réduire les temps d'attente.
          \item Éviter les conflits de ressources (bloc occupé, lit indisponible).
          \item Répartir la charge des équipes sur la semaine.
          \item Améliorer l'utilisation des ressources et réduire les coûts.
        \end{itemize}
      \end{block}

    \end{column}

  \end{columns}

  \vspace{0.4em}

  \begin{alertblock}{Problématique}
    Comment planifier les interventions pour optimiser les blocs opératoires, équilibrer l'occupation des lits et intégrer les urgences?
  \end{alertblock}

\end{frame}